\documentclass[10pt]{article}
\usepackage[margin=0.75in]{geometry}
\usepackage{amsmath,amssymb}
\usepackage{enumitem}
\usepackage{xcolor}
\usepackage{colortbl}
\usepackage{titlesec}
\usepackage{booktabs}
\usepackage{array}
\usepackage{multirow}
\usepackage{longtable}
\usepackage{tabularx}

\setlength{\parindent}{0pt}
\setlength{\parskip}{6pt}

% Colors
\definecolor{mustknow}{HTML}{FF4444}
\definecolor{highpri}{HTML}{FF9900}
\definecolor{medpri}{HTML}{FFCC00}
\definecolor{lowpri}{HTML}{88CC88}
\definecolor{tested}{HTML}{2266CC}
\definecolor{nottested}{HTML}{DDDDDD}
\definecolor{sectionbg}{HTML}{F0F4FF}

\titleformat{\section}{\large\bfseries}{}{0pt}{}[\vspace{-4pt}\rule{\textwidth}{0.5pt}]
\titleformat{\subsection}{\normalsize\bfseries}{}{0pt}{}

\newcommand{\YES}{\cellcolor{tested!20}\textbf{\textcolor{tested}{\checkmark}}}
\newcommand{\NO}{\cellcolor{nottested}}

\pagestyle{empty}

\begin{document}

{\centering\LARGE\bfseries ECO 310 Midterm 1 --- Preparation Strategy\par}
\vspace{2pt}
{\centering March 5, 2026 $\cdot$ In-class $\cdot$ 1 page cheat sheet (front \& back)\par}
\vspace{8pt}

% ============================================
\section{Topic Frequency Table (2022--2025 Practice Exams)}
% ============================================

Each checkmark means the topic appeared as a graded question (MC or short answer) on that exam. Priority is based on frequency across all four exams.

\vspace{4pt}

\renewcommand{\arraystretch}{1.15}
\begin{tabularx}{\textwidth}{>{\raggedright\arraybackslash}p{4.2cm} c c c c >{\centering\arraybackslash}p{1cm} >{\raggedright\arraybackslash}X}
\toprule
\textbf{Topic / Concept} & \textbf{'22} & \textbf{'23} & \textbf{'24} & \textbf{'25} & \textbf{Freq} & \textbf{Priority} \\
\midrule
\rowcolor{mustknow!15}
Marshallian demand derivation & \YES & \YES & \YES & \YES & 4/4 & \textcolor{mustknow}{\textbf{MUST-KNOW}} \\
\rowcolor{mustknow!15}
EU expression setup & \YES & \YES & \YES & \YES & 4/4 & \textcolor{mustknow}{\textbf{MUST-KNOW}} \\
\rowcolor{mustknow!15}
EU optimization via FOC & \YES & \YES & \YES & \YES & 4/4 & \textcolor{mustknow}{\textbf{MUST-KNOW}} \\
\rowcolor{mustknow!15}
Risk aversion (definition/test) & \YES & \YES & \YES & \YES & 4/4 & \textcolor{mustknow}{\textbf{MUST-KNOW}} \\
\rowcolor{mustknow!15}
Cobb-Douglas magic rule & \YES & \YES & \YES & \YES & 4/4 & \textcolor{mustknow}{\textbf{MUST-KNOW}} \\
\rowcolor{mustknow!15}
Perfect complements demand & \YES & \YES & \YES & \YES & 4/4 & \textcolor{mustknow}{\textbf{MUST-KNOW}} \\
\midrule
\rowcolor{highpri!15}
WARP / Revealed preference & \YES & \YES & \YES & \YES & 4/4 & \textcolor{highpri}{\textbf{HIGH}} \\
\rowcolor{highpri!15}
Sub/Income effect decomposition & \NO & \YES & \YES & \YES & 3/4 & \textcolor{highpri}{\textbf{HIGH}} \\
\rowcolor{highpri!15}
Certainty equiv.\ / Risk premium & \NO & \YES & \YES & \YES & 3/4 & \textcolor{highpri}{\textbf{HIGH}} \\
\rowcolor{highpri!15}
Corner solutions / non-negativity & \YES & \YES & \YES & \YES & 4/4 & \textcolor{highpri}{\textbf{HIGH}} \\
\rowcolor{highpri!15}
Convexity of preferences (MRS) & \YES & \YES & \YES & \NO & 3/4 & \textcolor{highpri}{\textbf{HIGH}} \\
\rowcolor{highpri!15}
Expenditure fn / Hicksian demand & \NO & \NO & \YES & \YES & 2/4 & \textcolor{highpri}{\textbf{HIGH}} \\
\rowcolor{highpri!15}
Portfolio allocation (risky+safe) & \NO & \NO & \NO & \YES & 1/4 & \textcolor{highpri}{\textbf{HIGH}} \\
\rowcolor{highpri!15}
Rationality (complete + transitive) & \YES & \YES & \YES & \NO & 3/4 & \textcolor{highpri}{\textbf{HIGH}} \\
\midrule
\rowcolor{medpri!15}
Perfect substitutes demand & \YES & \YES & \NO & \YES & 3/4 & \textcolor{medpri!80!black}{\textbf{MEDIUM}} \\
\rowcolor{medpri!15}
CES utility / limits & \NO & \NO & \NO & \YES & 1/4 & \textcolor{medpri!80!black}{\textbf{MEDIUM}} \\
\rowcolor{medpri!15}
Arrow-Pratt measures (ARA/RRA) & \NO & \NO & \NO & \YES & 1/4 & \textcolor{medpri!80!black}{\textbf{MEDIUM}} \\
\rowcolor{medpri!15}
Stochastic dominance (FOSD/SOSD) & \NO & \NO & \NO & \NO & 0/4 & \textcolor{medpri!80!black}{\textbf{MEDIUM}} \\
\rowcolor{medpri!15}
Edgeworth box / Gen.\ equilibrium & \NO & \NO & \NO & \NO & 0/4 & \textcolor{medpri!80!black}{\textbf{MEDIUM}} \\
\rowcolor{medpri!15}
Quasilinear utility & \YES & \NO & \NO & \NO & 1/4 & \textcolor{medpri!80!black}{\textbf{MEDIUM}} \\
\bottomrule
\end{tabularx}

\vspace{6pt}

\textbf{Repeated questions across exams:}
\begin{itemize}[nosep,leftmargin=12pt]
\item Apple tax + \$5 subsidy (Slutsky compensation) --- appeared in \textbf{both 2023 and 2024}
\item Framing paradox (consequentialism) --- appeared in \textbf{both 2023 and 2024}
\item 50-50 \$0/\$200 gamble preference --- appeared in \textbf{both 2023 and 2024}
\item Perfect complements sub/income effect --- 2023 (sub effect = 0) and 2024 (income effect = $-3$)
\end{itemize}

% ============================================
\section{Study Plan: Cheat Sheet + Practice Exams}
% ============================================

\subsection{Phase 1: Cheat Sheet Familiarity}
\begin{enumerate}[nosep,leftmargin=16pt]
\item \textbf{Read through the entire cheat sheet} end to end. For each yellow-highlighted formula, write it from memory on scratch paper, then check.
\item \textbf{Build a mental map:} Know which section contains which formula. The exam is timed---you cannot afford to search.
\item \textbf{Drill the recipes:} Section 9 of the cheat sheet has step-by-step recipes for demand derivation and EU problems. Practice following each recipe on a blank page until the steps are automatic.
\end{enumerate}

\subsection{Phase 2: Practice Exam Drilling}
\begin{enumerate}[nosep,leftmargin=16pt]
\item \textbf{Do 2025 exam first} (most recent, likely closest to your exam's style). Time yourself---treat it as real. Use only your cheat sheet.
\item \textbf{Then do 2024, 2023, 2022} in reverse chronological order.
\item After each exam, \textbf{grade yourself honestly}. For every mistake, identify:
  \begin{itemize}[nosep]
  \item Was it a \textbf{knowledge gap}? (Formula/concept not on cheat sheet or not understood)
  \item Was it a \textbf{procedure error}? (Forgot to check corners, wrong inequality direction, etc.)
  \item Was it a \textbf{reading error}? (Misidentified states, confused ``spend'' with ``buy'', etc.)
  \end{itemize}
\item \textbf{Update your cheat sheet} if you discover a gap.
\end{enumerate}

\subsection{Phase 3: Targeted Weakness Drilling}
\begin{enumerate}[nosep,leftmargin=16pt]
\item Collect every problem you got wrong or were slow on. Re-do them from scratch.
\item For each MUST-KNOW topic, you should be able to set up and solve a problem in under 3 minutes for MC, under 10 minutes for short answer.
\item Practice the EU optimization problems heavily---they appear on every exam and are the most algebra-intensive.
\end{enumerate}

% ============================================
\section{Problem Patterns to Drill}
% ============================================

\renewcommand{\arraystretch}{1.2}
\begin{tabularx}{\textwidth}{>{\raggedright\arraybackslash}p{3.5cm} >{\raggedright\arraybackslash}X >{\raggedright\arraybackslash}p{4cm}}
\toprule
\textbf{Pattern} & \textbf{What You Must Be Able To Do} & \textbf{Practice Sources} \\
\midrule
\textbf{Cobb-Douglas demand} & Apply magic rule instantly. Normalize exponents. State spending shares. & 2023 MC2, 2024 MC3, 2025 MC1 \\
\textbf{Perfect complements} & Set min args equal. Plug into budget. Know sub effect = 0. & 2023 MC4, 2024 MC5, 2025 MC3 \\
\textbf{Non-standard utility} demand & Compute MRS from scratch. Set MRS = price ratio + budget. Check corners. & 2022 Q3, 2023 Q3 \\
\textbf{WARP check} & Write cost comparisons at both price vectors. Determine violation range. & 2024 MC9, 2025 MC2, HW1 Prob 1 \\
\textbf{EU setup + FOC} & ID states/probs/wealth. Write EU. Differentiate. Solve. Check corners. & 2022 Q4, 2023 Q4, 2024 Q2b, 2025 SA2--SA3 \\
\textbf{CE / Risk premium} & Compute $EU$, invert $u$ to get CE, then $RP = E[W] - CE$. & 2024 Q2a, 2025 SA2a \\
\textbf{Sub/Income decomp} & Find original, compensated, and new bundles. Sub = comp $-$ orig. Inc = new $-$ comp. & 2023 MC4, 2024 MC5, 2025 MC6 \\
\textbf{Hicksian demand / EMP} & Minimize expenditure s.t.\ utility $\geq \bar{u}$. Not the same as Marshallian. & 2024 Q3d, 2025 MC3 \\
\textbf{Proof-style questions} & Risk aversion comparisons, DARA implications, EU inequality proofs. & 2023 Q2a, 2025 SA2a-ii, SA2b \\
\bottomrule
\end{tabularx}

% ============================================
\section{Common Exam Traps and Mistakes to Avoid}
% ============================================

\renewcommand{\arraystretch}{1.25}
\begin{tabularx}{\textwidth}{>{\centering\arraybackslash}p{0.4cm} >{\raggedright\arraybackslash}p{4cm} >{\raggedright\arraybackslash}X}
\toprule
& \textbf{Trap} & \textbf{How to Avoid It} \\
\midrule
1 & \textbf{Marshallian $\neq$ Hicksian} & Read the question. ``Cheapest bundle to reach utility $\bar{u}$'' = EMP/Hicksian. ``Best bundle given income $m$'' = UMP/Marshallian. The 2024 solutions note this cost many students half credit. \\
2 & \textbf{Perfect comp: sub effect = 0} & With $\min\{\}$, the compensated bundle equals the original bundle. The entire price-change effect is income effect. Do not compute a ``substitution effect.'' \\
3 & \textbf{WARP inequality directions} & ``$A$ revealed preferred to $B$'' means $A$ was chosen \emph{and} $B$ was affordable (cost of $B \leq$ cost of $A$ at those prices). Draw a table of costs. \\
4 & \textbf{EU: wrong wealth per state} & Use a probability tree. Final wealth = initial wealth $\pm$ all gains/losses in that state. Write it out explicitly before computing EU. \\
5 & \textbf{Forgetting non-negativity} & After solving the FOC, always check $x_i \geq 0$. If negative, set that good to 0 and re-solve. Especially common with quasilinear and perfect substitutes. \\
6 & \textbf{``Spend $k\times$'' $\neq$ ``Buy $k\times$''} & Cobb-Douglas magic rule gives \emph{spending} shares. ``Spends 3 times as much on good 2'' $\neq$ ``buys 3 times as many units'' unless prices are equal. Tested explicitly in 2024 MC3. \\
7 & \textbf{CD exponents not normalized} & $u = x^a y^b$: spending share on $x$ is $\frac{a}{a+b}$, NOT just $a$. Always divide by sum of exponents. \\
8 & \textbf{vNM: only affine transforms} & Expected utility rankings are preserved only by $\tilde{u} = a + bu$ ($b > 0$). General monotone transforms change EU rankings. Ordinal utility allows any monotone transform; cardinal (vNM) does not. \\
9 & \textbf{Variance $\neq$ Std deviation} & Portfolio std dev $= \alpha \cdot \sigma_{\text{risky}}$. If problem gives variance, take the square root first. Tested in 2025 MC4. \\
10 & \textbf{Quasilinear: income effect for $x_1$} & At interior, demand for $x_1$ is independent of $m$. But if $m$ is too low, you hit a corner at $x_2 = 0$, and then $x_1 = m/p_1$ \emph{does} depend on income. \\
11 & \textbf{Risk aversion $\neq$ always choose certain} & A risk-averse person prefers \$100 certain over 50-50 \$0/\$200 (same EV). But \$90 certain vs.\ 50-50 \$0/\$200 is \textbf{ambiguous}---lower EV but no risk. Cannot determine without knowing the utility function. \\
12 & \textbf{CES limit behavior} & As $r \to 0$: Cobb-Douglas. As $r \to -\infty$: perfect complements. As $r = 1$: perfect substitutes. Know which limit is which. \\
\bottomrule
\end{tabularx}

\vspace{8pt}

\begin{center}
\fbox{\parbox{0.9\textwidth}{\centering\textbf{Bottom line:} Every exam has (1) a Cobb-Douglas/complements demand MC, (2) a WARP question, (3) a non-standard demand derivation, and (4) an EU optimization problem. Master these four patterns and you cover $>$80\% of the exam.}}
\end{center}

\end{document}
